\documentclass[11pt]{article}
\usepackage[a4paper,margin=26mm]{geometry}
\usepackage{amsmath,amssymb,amsthm,booktabs,hyperref}
\newtheorem{theorem}{Theorem}
\newtheorem{lemma}[theorem]{Lemma}
\title{Two exact physical accelerations for PA194AX11GH and PA194AX11GHS}
\author{Research supplement}
\date{14 November 2026}
\begin{document}
\maketitle
\noindent\textbf{Status.} This supplement proves a relative machine transformation.
Absolute calibration to PA consistency is treated separately in the enclosing
report and its cited incoming written proof.
The proof below is an ordinary mathematical proof, explicitly using only
primitive recursion and quantifier-free induction for its arithmetization;
it is not described as a compiled Lean theorem.

\section{The two changes}
Native controls and symbols are used throughout. The initial control is 74,
the halt label is 254, and the tape alphabet is $\{0,1\}$ with blank 0.
Let $M$ be either PA194AX11GH or PA194AX11GHS. Define $M^{\mathrm F}$ by
changing exactly the following two transitions and retaining every other row:
\begin{center}
\begin{tabular}{ccccc}
\toprule
Control & Read & Original action & Accelerated action & Old steps / new steps\\
\midrule
26 & 1 & $1R42$ & $1R91$ & $3/1$\\
101 & 0 & $0L26$ & $0R17$ & $3/1$\\
\bottomrule
\end{tabular}
\end{center}
Here $wDq$ means write $w$, move in direction $D$, and enter control $q$.
Both resulting tables still have 194 working controls and 388 transitions.
The dense versions start at 0 and halt at label 194; the accompanying label
arrays give the exact conjugacy with native numbering.

\section{A finite structural invariant}
Write a configuration as $(q,h,T)$, where $T(i)$ is the symbol at position $i$.
Consider the following three implications:
\begin{align*}
I_{101}:\quad &q=101\ \Longrightarrow\ T(h-1)=0,\\
I_{11}:\quad &q=11\ \Longrightarrow\bigl(T(h-1)=0\ \lor\ T(h+1)=0\bigr),\\
I_{26}:\quad &q=26\ \land\ T(h)=1\ \Longrightarrow\ T(h+1)=T(h+2)=0.
\end{align*}
Let $I$ be their conjunction. The finite row checker establishes:
\begin{enumerate}
\item The sole transition into 101 is $96/0\mapsto0R101$.
\item Every transition into 11 writes 0. Their source pairs are
$(14,0),(39,0),(56,0),(312,0)$; the first three move left and the last moves right.
\item The only transitions into 26 are $11/0\mapsto0L26$ and
$101/0\mapsto0L26$.
\end{enumerate}
The start control 74 is different from 11, 26 and 101.

\begin{lemma}
Every parent run beginning in control 74 satisfies $I$ at every time,
on every tape and at every initial head position.
\end{lemma}
\begin{proof}
The invariant is initially vacuous. A step into 101 has just written 0 to
the cell immediately left of the new head, proving $I_{101}$. A step into 11
has written 0 to one of its two neighboring cells, proving $I_{11}$.

For a step into 26, both possible incoming transitions move left after
writing 0, so the new cell immediately to the right is 0. If the predecessor
control is 101, $I_{101}$ makes the newly scanned symbol 0, and the antecedent
of $I_{26}$ is false. If the predecessor is 11 and the new scanned symbol is 1,
$I_{11}$ forces the old right neighbor to be 0. This is the new cell two
positions to the right. Thus $I_{26}$ holds. No other step enters the three
relevant controls, completing the induction.
\end{proof}

This argument does not invoke source-program correctness, a register-file
layout, a sampled execution, or an arithmetic-box contract.

\section{Exact local identities}
At control 26 reading 1, $I_{26}$ supplies the next two symbols $00$.
The original sequence is
\[
 (26,h,T)\ \xrightarrow{1R42}\ (42,h+1,T)
 \ \xrightarrow{0R13}\ (13,h+2,T)
 \ \xrightarrow{0L91}\ (91,h+1,T).
\]
All writes preserve the scanned symbols. This is exactly the configuration
reached in one accelerated step $1R91$.

At control 101 reading 0, $I_{101}$ supplies a 0 immediately to the left.
The original sequence is
\[
 (101,h,T)\ \xrightarrow{0L26}\ (26,h-1,T)
 \ \xrightarrow{0R194}\ (194,h,T)
 \ \xrightarrow{0R17}\ (17,h+1,T),
\]
again exactly the endpoint of the new single step $0R17$.
None of the displayed intermediate controls is the halt control.
The checker independently enumerates the four length-five windows for the
first guard and the eight length-five windows for the second, checking full
tape equality, head equality and control equality in all twelve cases.

\section{Uniform simulation and weak arithmetic}
For arithmetic coding, extend both step functions by making the halting
configuration absorbing. This convention is used only in the arithmetized
iteration; it does not add a working control to the machine.
Let $C_t^{\mathrm F}$ be the accelerated configuration after $t$ steps and set
\[
 \kappa(C)=\begin{cases}
 3,&(q,T(h))\in\{(26,1),(101,0)\},\\
 1,&\text{otherwise}.
 \end{cases}
\]
Define the primitive-recursive expanded clock by
\[
 \sigma(0)=0,\qquad
 \sigma(t+1)=\sigma(t)+\kappa(C_t^{\mathrm F}).
\]

\begin{theorem}
For a common initial configuration with control 74,
\[
 C_t^{\mathrm F}=C_{\sigma(t)},\qquad
 t\leq\sigma(t)\leq3t.
\]
For the primitive-recursive, absorbing-clock halting predicates $H_M$ and
$H_{M^{\mathrm F}}$, iPRA proves
\[
 \forall t\,\bigl(H_M(t)\ \longrightarrow\ H_{M^{\mathrm F}}(t)\bigr)
 \quad\text{and}\quad
 \forall t\,\bigl(H_{M^{\mathrm F}}(t)\ \longrightarrow\ H_M(3t)\bigr).
\]
Consequently iPRA proves the relative blank-start halting equivalence and
the corresponding nonhalting equivalence.
\end{theorem}
\begin{proof}
Induct on $t$. A configuration satisfying $I$ has its accelerated successor
equal to the parent successor after either one or three steps, by the exact
identities above. The parent invariant is preserved throughout those steps.
The clock equation and its bounds follow by quantifier-free induction.

If the parent has halted by $t$, then its absorbing state at $\sigma(t)\geq t$
is also halted, so the accelerated configuration at $t$ is halted. Conversely,
if the accelerated configuration at $t$ is halted, so is the parent at
$\sigma(t)\leq3t$ and hence at $3t$. Both implications are constructive.
The configuration functions, the finite invariant, the case selector
$\kappa$, and $\sigma$ are primitive recursive on the usual finite-support
tape codes. Their displayed induction predicates are quantifier-free.
Thus the same proof is available in iPRA, in particular for the blank tape.
No soundness or consistency principle for PA is used.
\end{proof}

When either machine halts, its least halting times obey
$\tau_{M^{\mathrm F}}\leq\tau_M\leq3\tau_{M^{\mathrm F}}$.
The result also gives
\[
 \operatorname{Con}(\mathrm{PA})\leftrightarrow\neg\operatorname{Halt}(M)
 \quad\Longrightarrow\quad
 \operatorname{Con}(\mathrm{PA})\leftrightarrow\neg\operatorname{Halt}(M^{\mathrm F})
\]
as an implication provable in iPRA. Its application requires the separate
parent calibration argument. That argument's mathematical and formalization
status is assessed in the enclosing report; it is not part of the shortcut proof.

\section{Reproduction and evidence scope}
Run \texttt{python3 acceleration/build\_fast.py --full} from the enclosing
AX11GH/GHS packet. The script checks every structural hypothesis above,
the twelve local identities, complete source contracts, dense/native
conjugacy, and whole blank initialization. During initialization it checks
the exact expanded-clock identity at every accelerated transition and
compares the entire finite tape support. A separate literal checker derives
the list-sweep lemmas from the rows and repeats all source outcomes without
importing the construction compiler or arithmetic-box oracle.

All exact counts, file hashes and shortcut-use counts are recorded in the
per-machine JSON receipts. The transition tables and receipts are the
machine-specific evidence; the invariant and simulation theorem above
explain why the acceleration preserves the parent's consistency statement.
\end{document}
